\documentclass[12pt]{article}
\makeatletter
\def\input@path{{../../}}
\makeatother
\usepackage{sbc-template}
\makeatletter
\renewcommand{\inst}[1]{}
\makeatother
\usepackage{graphicx,url}
\usepackage[utf8]{inputenc}
\usepackage[T1]{fontenc}
\usepackage[brazil]{babel}
\usepackage{longtable,booktabs,array}
\providecommand{\tightlist}{%%
  \setlength{\itemsep}{0pt}\setlength{\parskip}{0pt}}
\sloppy
\title{Quantum Reservoir Computing para previsao de demanda: um estudo didatico com o Favorita}
\author{}
\address{}
\begin{document}

\maketitle

\begin{resumo}
Este artigo fecha a serie com um estudo aplicado de Quantum Reservoir Computing (QRC) no mesmo problema de negocio introduzido no primeiro texto: previsao de demanda e apoio a decisoes de estoque. Usando o recorte \texttt{store\_nbr\ =\ 1}, \texttt{family\ =\ BEVERAGES} do dataset Favorita e o mesmo protocolo temporal dos artigos anteriores, implementamos um pipeline de QRC em Qiskit, com codificacao sequencial de entrada, dinamica quantica fixa, observaveis \texttt{Z}, \texttt{X} e \texttt{ZZ}, e readout Ridge classico. O melhor ponto por criterio multi-metrica ficou em \texttt{14} qubits com \texttt{window\ =\ 7}, produzindo \texttt{MAE\ =\ 441.286}, \texttt{RMSE\ =\ 525.325}, \texttt{WAPE\ =\ 0.2038} e \texttt{sMAPE\ =\ 0.2269}. Uma funcao objetivo escalar com penalizacao leve de complexidade, por sua vez, preferiu \texttt{8} qubits com \texttt{window\ =\ 5}, destacando o compromisso entre qualidade e custo. Embora o QRC nao supere RC nem as tecnicas classicas neste recorte, o estudo cumpre um papel central: mostrar como construir, avaliar e interpretar um experimento inicial de QRC em um problema realista sem perder honestidade metodologica.
\end{resumo}

\section{1. Introducao}\label{1-introducao}

A pergunta que abre o artigo final da serie nao e "o QRC vence?". A pergunta correta e "como devemos montar e ler um experimento de QRC em um problema real?".

Se o leitor percorreu os seis artigos anteriores, ja sabe que:

\begin{itemize}
\tightlist
\item
  o problema de negocio e previsao de demanda;
\item
  o benchmark classico esta estabelecido;
\item
  Reservoir Computing e um paradigma de dinamica rica mais readout simples;
\item
  Quantum Reservoir Computing herda essa estrutura, trocando o suporte dinamico.
\end{itemize}

O objetivo agora e fechar a obra com um estudo completo, pequeno o suficiente para ser reproduzido e serio o suficiente para ser comparado.

\section{2. Recorte experimental herdado da serie}\label{2-recorte-experimental-herdado-da-serie}

O experimento final preserva o mesmo recorte central:

\begin{itemize}
\tightlist
\item
  dataset: Favorita \texttt{train.csv}
\item
  serie: \texttt{store\_nbr\ =\ 1}
\item
  familia: \texttt{BEVERAGES}
\item
  frequencia: diaria
\item
  teste: ultimos 90 dias
\end{itemize}

Essa escolha importa porque evita um vicio comum em estudos emergentes: mudar o problema no momento de introduzir o metodo mais novo.

Aqui, o problema e o mesmo. O que muda e o tipo de reservatorio.

\section{3. Pipeline de QRC usado no projeto}\label{3-pipeline-de-qrc-usado-no-projeto}

No repositorio, o QRC esta implementado em \texttt{code/qrc/model.py}. O pipeline segue uma versao minima e didatica do paradigma.

\subsection{3.1 Entrada}\label{31-entrada}

O vetor quantico de entrada usa quatro sinais derivados do pipeline sequencial:

\begin{itemize}
\tightlist
\item
  demanda previa normalizada;
\item
  promocao normalizada;
\item
  \texttt{dow\_sin};
\item
  \texttt{dow\_cos}.
\end{itemize}

Esses sinais sao reencodificados ao longo da janela temporal.

\subsection{3.2 Dinamica quantica}\label{32-dinamica-quantica}

Cada passo temporal aplica portas parametrizadas \texttt{RX} e \texttt{RY} em todos os qubits, seguidas por uma cadeia de acoplamentos \texttt{CZ} com rotações \texttt{RZ} em anel. Essa dinamica e fixa apos a inicializacao aleatoria controlada por \texttt{seed}.

\subsection{3.3 Observaveis}\label{33-observaveis}

O readout nao observa diretamente um "estado interno" classico. Ele recebe valores esperados de:

\begin{itemize}
\tightlist
\item
  observaveis \texttt{Z} por qubit;
\item
  observaveis \texttt{X} por qubit;
\item
  observaveis \texttt{ZZ} entre vizinhos.
\end{itemize}

\subsection{3.4 Readout}\label{34-readout}

Assim como no RC classico da serie, o readout continua sendo simples: uma regressao Ridge treinada sobre os observaveis.

Essa continuidade e intencional. Ela faz com que a comparacao destaque o reservatorio, e nao uma troca de cabeca preditiva.

\section{4. Configuracao experimental}\label{4-configuracao-experimental}

O tuning do QRC foi conduzido em duas camadas:

\begin{enumerate}
\def\labelenumi{\arabic{enumi}.}
\tightlist
\item
  criterio puramente multi-metrica;
\item
  criterio escalar com penalizacao leve de complexidade.
\end{enumerate}

\subsection{4.1 Melhor configuracao por criterio multi-metrica}\label{41-melhor-configuracao-por-criterio-multi-metrica}

Depois de extrapolar a busca ate \texttt{16} qubits e janelas maiores, a configuracao balanceada mais forte ficou em:

\begin{itemize}
\tightlist
\item
  \texttt{n\_qubits\ =\ 14}
\item
  \texttt{window\ =\ 7}
\item
  \texttt{input\_scale\ =\ 1.2}
\item
  \texttt{ridge\_alpha\ =\ 0.1}
\end{itemize}

\subsection{4.2 Melhor configuracao por funcao objetivo escalar}\label{42-melhor-configuracao-por-funcao-objetivo-escalar}

Tambem foi implementada uma funcao objetivo:

\texttt{J(q,\ w)\ =\ G(q,\ w)\ *\ C(q,\ w)}

onde \texttt{G} e a media geometrica ponderada das razoes entre \texttt{MAE}, \texttt{RMSE}, \texttt{WAPE}, \texttt{sMAPE} do QRC e as metricas de referencia do benchmark nao-QRC, enquanto \texttt{C} adiciona uma penalizacao leve por complexidade em \texttt{qubits} e \texttt{window}.

Sob esse criterio, a configuracao preferida no conjunto focal ficou em:

\begin{itemize}
\tightlist
\item
  \texttt{n\_qubits\ =\ 8}
\item
  \texttt{window\ =\ 5}
\end{itemize}

Essa distincao e importante. \texttt{14x7} aperta melhor as metricas do QRC. \texttt{8x5} oferece melhor compromisso entre qualidade e custo de complexidade.

\section{5. Resultados do QRC}\label{5-resultados-do-qrc}

\subsection{5.1 Resultado do QRC 14x7}\label{51-resultado-do-qrc-14x7}

No recorte principal, o QRC \texttt{14x7} obteve:

\begin{longtable}[]{@{}lrrrr@{}}
\toprule\noalign{}
Configuracao & MAE & RMSE & WAPE & sMAPE \\
\midrule\noalign{}
\endhead
\bottomrule\noalign{}
\endlastfoot
QRC 14x7 & 441.286 & 525.325 & 0.2038 & 0.2269 \\
\end{longtable}

\subsection{5.2 Resultado do QRC 8x5}\label{52-resultado-do-qrc-8x5}

Rodando o mesmo protocolo com a configuracao favorecida pela funcao objetivo escalar:

\begin{longtable}[]{@{}lrrrr@{}}
\toprule\noalign{}
Configuracao & MAE & RMSE & WAPE & sMAPE \\
\midrule\noalign{}
\endhead
\bottomrule\noalign{}
\endlastfoot
QRC 8x5 & 455.833 & 526.949 & 0.2105 & 0.2342 \\
\end{longtable}

\subsection{5.3 Leitura direta}\label{53-leitura-direta}

No run comparativo com \texttt{seed\ =\ 7}, \texttt{14x7} ficou melhor do que \texttt{8x5} nas quatro metricas. Ainda assim, \texttt{8x5} continua sendo a escolha mais economica segundo a funcao objetivo penalizada por complexidade.

\section{6. Comparacao com RC e tecnicas classicas}\label{6-comparacao-com-rc-e-tecnicas-classicas}

O quadro abaixo posiciona o QRC dentro do benchmark da serie.

\begin{longtable}[]{@{}llrrrr@{}}
\toprule\noalign{}
Modelo & Familia & MAE & RMSE & WAPE & sMAPE \\
\midrule\noalign{}
\endhead
\bottomrule\noalign{}
\endlastfoot
ETS & baseline estatistico & 216.303 & 307.638 & 0.0999 & 0.1042 \\
XGBoost & ML classico & 256.666 & 332.828 & 0.1185 & 0.1237 \\
Prophet & baseline estatistico & 257.584 & 321.818 & 0.1189 & 0.1319 \\
Seasonal Naive & baseline estatistico & 259.067 & 352.280 & 0.1196 & 0.1361 \\
RC & reservoir computing classico & 324.294 & 423.417 & 0.1498 & 0.1653 \\
LSTM & DL classico & 371.827 & 461.342 & 0.1717 & 0.1822 \\
QRC 14x7 & reservoir computing quantico & 441.286 & 525.325 & 0.2038 & 0.2269 \\
QRC 8x5 & reservoir computing quantico & 455.833 & 526.949 & 0.2105 & 0.2342 \\
\end{longtable}

Essa tabela permite tres conclusoes claras.

\subsection{6.1 O QRC funciona de ponta a ponta}\label{61-o-qrc-funciona-de-ponta-a-ponta}

Isso pode parecer pouco ambicioso, mas nao e. Em uma serie introdutoria, mostrar um pipeline quantico funcional e comparavel e um resultado importante.

\subsection{6.2 O QRC ainda nao vence os classicos neste recorte}\label{62-o-qrc-ainda-nao-vence-os-classicos-neste-recorte}

ETS, XGBoost, Prophet e ate o Seasonal Naive continuam fortes. O RC classico tambem supera o QRC atual.

\subsection{6.3 A interpretacao correta nao e "fracasso"}\label{63-a-interpretacao-correta-nao-e-fracasso}

O valor do experimento esta em estabelecer:

\begin{itemize}
\tightlist
\item
  um protocolo comparavel;
\item
  um pipeline quantico claro;
\item
  um criterio de tuning explicito;
\item
  uma leitura honesta de limites e potencial.
\end{itemize}

\section{7. O que os resultados mostram e o que eles nao mostram}\label{7-o-que-os-resultados-mostram-e-o-que-eles-nao-mostram}

\subsection{7.1 O que mostram}\label{71-o-que-mostram}

\begin{itemize}
\tightlist
\item
  QRC pode ser implementado no mesmo problema de negocio usado pelos classicos;
\item
  tuning de \texttt{qubits} e \texttt{window} altera significativamente o desempenho;
\item
  criterios diferentes de otimizacao produzem escolhas diferentes;
\item
  aumentar qubits nao garante melhora robusta.
\end{itemize}

\subsection{7.2 O que nao mostram}\label{72-o-que-nao-mostram}

\begin{itemize}
\tightlist
\item
  que QRC seja superior a baselines fortes neste problema;
\item
  que a configuracao atual seja otima globalmente;
\item
  que o comportamento observado generalize para todo o Favorita;
\item
  que simulacao statevector equivalha a execucao em hardware quantico real.
\end{itemize}

Explicitar essas fronteiras e essencial para fechar a serie com rigor.

\section{8. A contribuicao didatica do estudo final}\label{8-a-contribuicao-didatica-do-estudo-final}

Mesmo ficando abaixo dos modelos classicos neste recorte, o artigo final entrega uma contribuicao pedagogica forte:

\begin{enumerate}
\def\labelenumi{\arabic{enumi}.}
\tightlist
\item
  traduz QRC em um pipeline concreto;
\item
  mostra como comparar QRC sem privilégio metodologico;
\item
  introduz uma funcao objetivo clara para tuning;
\item
  ensina o leitor a distinguir desempenho puro de desempenho ajustado por complexidade.
\end{enumerate}

Isso faz do artigo final nao apenas um estudo aplicado, mas tambem uma sintese metodologica da obra inteira.

\section{9. O que o leitor aprendeu do artigo 1 ao 7}\label{9-o-que-o-leitor-aprendeu-do-artigo-1-ao-7}

Ao fim da serie, o leitor foi conduzido por um caminho continuo:

\begin{enumerate}
\def\labelenumi{\arabic{enumi}.}
\tightlist
\item
  entender o problema de negocio de previsao de demanda;
\item
  formular o caso Favorita como problema temporal reproduzivel;
\item
  aprender o paradigma de RC;
\item
  construir um primeiro benchmark com baselines;
\item
  comparar RC com IA classica;
\item
  entender a ponte entre RC classico, reservatorios fisicos e QRC;
\item
  montar e interpretar um experimento inicial de QRC.
\end{enumerate}

Em outras palavras, a obra cumpre o papel de uma introducao 101 a QRC partindo do zero e chegando a um estudo aplicado.

\section{10. Conclusao}\label{10-conclusao}

O resultado final desta serie nao e um trofeu de benchmark para QRC. E algo mais solido: um caminho didatico completo, tecnicamente honesto e experimentalmente reproduzivel para ensinar Quantum Reservoir Computing a partir de um problema real de negocio.

Se o leitor sair desta obra entendendo por que QRC e uma extensao natural de RC, como montar um pipeline quantico minimo e como comparar esse pipeline com modelos classicos sem autoengano, entao a serie cumpriu seu objetivo.

\section{Entregaveis associados no repositorio}\label{entregaveis-associados-no-repositorio}

\begin{itemize}
\tightlist
\item
  implementacao do QRC: \texttt{code/qrc/model.py}
\item
  resultado principal do QRC: \texttt{code/qrc/RESULTS.md}
\item
  funcao objetivo: \texttt{code/qrc/OBJECTIVE.md}
\item
  tuning de qubits e janela: \texttt{code/qrc/experiments/TUNING.md}
\item
  benchmark consolidado: \texttt{code/RESULTS.md}
\end{itemize}

\section{Referencias}\label{referencias}

\begin{enumerate}
\def\labelenumi{\arabic{enumi}.}
\tightlist
\item
  Fujii, K., Nakajima, K. Harnessing disordered-ensemble quantum dynamics for machine learning.
\item
  Tanaka, G. et al. Recent advances in physical reservoir computing: A review.
\item
  Jaeger, H. The "echo state" approach to analysing and training recurrent neural networks.
\item
  Chen, T., Guestrin, C. XGBoost: A scalable tree boosting system.
\item
  Taylor, S. J., Letham, B. Forecasting at scale.
\end{enumerate}

\section*{Resultados Computacionais Associados}

Os resultados computacionais associados a este artigo estao organizados na subpasta \texttt{computational\_results\_20260402\_222902/}, localizada dentro desta pasta \LaTeX. Esse diretorio contem o arquivo \texttt{REPORT.md}, tabelas em CSV e imagens comparativas apropriadas ao escopo do artigo.

\end{document}
